% Status: Anhang mit ergaenzenden Abbildungen (4b).
\chapter{Anhang: Ergänzende Abbildungen}
\label{ch:anhang}

Dieser Anhang sammelt ergänzende, im Haupttext nicht abgebildete Diagnostiken zu H1
sowie eine Ansicht des Monitoring-Werkzeugs. Sie stützen die Befunde, sind für die
Argumentation aber nicht zwingend.

\section{H1-Robustheit nach Prognosehorizont}
Abbildung~\ref{fig:anhang-horizon} schlüsselt den H1-Poll-Vergleich nach
Prognosehorizont auf und zeigt, dass der Polymarket-Vorteil an das Fenster bis 90
Tage gebunden ist, während das volle Panel ihn nicht trägt.

\begin{figure}[ht]
\centering
\includegraphics[width=0.95\textwidth]{h1_state_poll_panel_horizon_diagnostic.png}
\caption{H1-Robustheit nach Prognosehorizont: mittlerer Brier-Verlust,
Lower-Loss-Zählung und Mehrheitsschwelle je Horizont-Bin sowie Claim-Audit. Das
Fenster bis 90 Tage stützt Polymarket, das volle Panel nicht. Quelle:
\texttt{data/results/h1\_state\_poll\_panel\_horizon\_diagnostic.png}.}
\label{fig:anhang-horizon}
\end{figure}

\section{H1-Robustheit gegenüber Ausreissern und Bundesstaaten}
Abbildung~\ref{fig:anhang-outlier} prüft, ob der Vorteil an einzelnen Bundesstaaten
hängt.

\begin{figure}[ht]
\centering
\includegraphics[width=0.95\textwidth]{h1_direct_poll_outlier_robustness.png}
\caption{H1-Robustheit gegenüber Ausreissern: Leave-one-state-out-Mittel, Erosion
beim Entfernen der grössten positiven Staaten und Beitrag je Bundesstaat. Der
Vorteil hängt nicht an einem einzelnen Staat, bleibt aber auf die Top-Staaten
konzentriert; die breite Mehr-Fälle-Behauptung bleibt ungedeckt. Quelle:
\texttt{data/results/h1\_direct\_poll\_outlier\_robustness.png}.}
\label{fig:anhang-outlier}
\end{figure}

\section{Monitoring-Werkzeug (Read-only-Prototyp)}
Abbildung~\ref{fig:anhang-monitor} zeigt eine Ansicht des read-only Monitors über
mehrere Live-Märkte.

\begin{figure}[ht]
\centering
\includegraphics[width=0.95\textwidth]{monitor_v2_polymarket_rolling_history.png}
\caption{Monitoring-Werkzeug (Read-only-Prototyp): rollende Polymarket-Midpoint-Historie
und aggregierte öffentliche Handelsaktivität je Zeitfenster (Illustration, kurzer
Beobachtungszeitraum). Quelle:
\texttt{data/results/monitor\_v2\_polymarket\_rolling\_history.png}.}
\label{fig:anhang-monitor}
\end{figure}
